\documentclass[11pt,letterpaper]{article}

\usepackage[margin=0.75in]{geometry}
\usepackage[T1]{fontenc}
\usepackage{lmodern}
\usepackage{microtype}
\usepackage{titlesec}
\usepackage[hidelinks]{hyperref}

\setlength{\parindent}{0pt}
\setlength{\parskip}{3pt}
\titleformat{\section}{\normalsize\bfseries}{}{0pt}{}
\titlespacing*{\section}{0pt}{6pt}{3pt}
\urlstyle{same}

\newcommand{\groupsize}{2}
\newcommand{\groupmembers}{Ethan Ku and Raghav Ramkumar}

\begin{document}

\begin{center}
    {\large\bfseries Predicting the Effect of Layout on Analog Integrated Circuit Performance}\\[5pt]
    Group size: \groupsize \qquad \groupmembers
\end{center}

\section*{Problem Statement}

Analog circuits can behave differently after layout because wires add resistance
and capacitance. These effects change the circuit's response to different signal
frequencies. Simulating every possible layout is time-consuming. We will test
whether machine learning can predict these changes and help choose promising
layouts. This is challenging because device placement, connections, and electrical
behavior affect one another~\cite{osiris}.

\section*{Methods}

Ridge regression will provide a linear baseline with regularization to reduce
overfitting. Kernel ridge regression with an RBF kernel will use similarity
between examples to fit nonlinear patterns. A small neural network will learn
interactions between features. These approaches build on regression, kernels,
and neural networks covered in class.

We will predict each of OSIRIS's three parasitic extraction (PEX) scores, which
summarize differences between simulated responses before and after layout.
We will first use device information, including transistor finger counts
(the number of parallel sections in a transistor). We will then add device
positions, distances, and symmetry measurements to test whether geometry
improves predictions.

\section*{Related Work}

OSIRIS supplies labeled layout variants and a layout-optimization
baseline~\cite{osiris}. ALIGN studies automatic layout generation from circuit
descriptions~\cite{align}. Dhar et al.\ use SVMs and neural networks to check
electrical constraints during layout design~\cite{dhar}. We will use existing
layouts to compare prediction methods and measure how well their scores help
select a layout.

\section*{Plan of Experiments}

The free OSIRIS repository on Hugging Face contains about 89,000 simulated layout
variants across five circuit types, including connection files (netlists), layout
geometry (GDS), and simulation results before and after layout~\cite{dataset}.
These paired results let us measure the effect of layout. We will use the
repository's loading example and released code to check data definitions, then
write our own scripts using Python, scikit-learn, PyTorch, and gdstk.
We will train separate models for each circuit type, starting with the Miller
operational transconductance amplifier (OTA).

We will use 70\% training, 15\% validation, and 15\% test data. Within each circuit
type, all layouts with the same transistor finger counts will stay in one set,
so related layouts cannot appear in both training and testing.
Scaling will use training data, and model settings will be chosen using
validation data. We will compare models with and without geometry against a
predictor that always returns the training mean. We will report mean absolute
error (MAE) and root mean square error (RMSE) for each score and circuit type.
For each test configuration, we will select the layout with the lowest predicted
score and measure the gap between its actual score and the best available score.

\section*{Project Plan}

Ethan and Raghav will both prepare the data and extract layout features. We will implement
and train the models together as well. Both will review results and write the report.
Our milestones are to (1) verify labels and data splits and obtain a ridge
baseline, (2) add geometry features and train the other models,
(3) compare prediction errors and layout selection, and (4) analyze errors and
finish the report.
Matching layout shapes to circuit devices may be difficult. If necessary, we
will focus on one OTA type and use simple component positions and distances.
If kernel training exceeds available memory, we will compare all models on
the same smaller training set, then expand ridge and neural-network training
to the full training set.

\newpage

\begin{thebibliography}{9}

\bibitem{osiris}
G.~Chiari, M.~Piccoli, and D.~Zoni.
OSIRIS: Bridging Analog Circuit Design and Machine Learning with Scalable
Dataset Generation.
\textit{International Conference on Learning Representations}, 2026.
\href{https://proceedings.iclr.cc/paper_files/paper/2026/hash/9c0be694049f5b6fb35ed9ef12da34b0-Abstract-Conference.html}{Conference paper}.

\bibitem{align}
K.~Kunal et al.
ALIGN: Open-Source Analog Layout Automation from the Ground Up.
\textit{Proceedings of the 56th Annual Design Automation Conference}, 2019.
\url{https://doi.org/10.1145/3316781.3323471}.

\bibitem{dhar}
T.~Dhar et al.
Fast and Efficient Constraint Evaluation of Analog Layout Using Machine
Learning Models.
\textit{Proceedings of the 26th Asia and South Pacific Design Automation
Conference}, pp.~158--163, 2021.
\url{https://doi.org/10.1145/3394885.3431547}.

\bibitem{dataset}
Hardware-Fab.
OSIRIS Dataset and Code.
Hugging Face, accessed November 5, 2026.
Dataset, download instructions, loading example, and pipeline code:
\url{https://huggingface.co/datasets/hardware-fab/osiris}.

\end{thebibliography}

\end{document}